\documentclass[aspectratio=169]{beamer}

% ==========================================================================
%  presentation_results_benchmarking.tex
%
%  Results / benchmarking half of the ShinySwarm deck. Companion to
%  presentation_intro_and_rqs.tex (same theme + colours).
%
%  Each results slide maps to one figure produced by analysis/make-plots.R:
%    fig01_relay_per_app        -> Slide: state relay latency
%    fig02_throughput_vs_load   -> Slide: latency vs concurrent users
%    fig03_throughput_per_app   -> Slide: throughput latency per app
%    fig04_save_restore         -> Slide: save / restore latency
%    fig05_overview_grid        -> Slide: latency overview across benchmarks
%    fig06_reliability          -> Slide: data loss / error / isolation
%    fig07_slowdown_factor      -> Slide: slowdown factor under load
%
%  Figures live in analysis/plots/. Until the real benchmark CSVs are
%  collected, each \includegraphics is wrapped so the deck still compiles
%  with a labelled placeholder instead of a missing-file error.
%
%  Build:
%    cd docs/thesis && pdflatex presentation_results_benchmarking.tex
%
%  Regenerate figures first (real or mock):
%    Rscript analysis/make-plots.R k6/results analysis/plots rest kafka
% ==========================================================================

% --- THEME (identical to the intro deck) -----------------------------------
\usetheme{Madrid}
\usecolortheme{seahorse}
\setbeamertemplate{navigation symbols}{}
\setbeamertemplate{footline}[frame number]
\setbeamercolor{frametitle}{bg=structure.fg!15!bg}

\usepackage[utf8]{inputenc}
\usepackage[T1]{fontenc}
\usepackage{booktabs}
\usepackage{graphicx}
\usepackage{tikz}
\usetikzlibrary{shapes.geometric, arrows.meta, positioning, fit, calc}
\usepackage{hyperref}
\usepackage{fontawesome5}

% --- figure directory + safe-include helper --------------------------------
% Resolves whether built from docs/thesis/ or from the repo root.
\graphicspath{{../../analysis/plots/}{analysis/plots/}{./}}

% \safefig{<basename>}{<width>}  -- include the figure if present, else a
% labelled placeholder box so the slide never breaks the build.
\newcommand{\safefig}[2]{%
  \IfFileExists{../../analysis/plots/#1.png}%
    {\includegraphics[width=#2]{#1.png}}%
    {\IfFileExists{analysis/plots/#1.png}%
      {\includegraphics[width=#2]{#1.png}}%
      {\fbox{\parbox[c][3.2cm][c]{#2}{\centering\footnotesize\ttfamily
        \detokenize{#1}.png\\[0.4em]\normalfont\itshape
        figure not yet generated\\\scriptsize
        run \texttt{make-plots.R} once benchmark CSVs exist}}}}%
}

\title{\textbf{ShinySwarm} --- Benchmarking Results}
\subtitle{Latency, Throughput, Reliability:\\Monolithic vs.\ REST vs.\ Kafka}
\author{Iva Gunjaca}
\institute{Master of Software Engineering\\University of Amsterdam\\[0.5em]
\small Supervisor: Dr Zhiming Zhao\\
MultiScale Networked Systems Group, Informatics Institute}
\date{Results Presentation --- 2026}

\begin{document}

% ============================================================
% SLIDE 1: TITLE
% ============================================================
\begin{frame}
\titlepage
\end{frame}

% ============================================================
% SLIDE 2: ROADMAP
% ============================================================
\begin{frame}{From Research Question to Evidence}

\begin{columns}[T]
\begin{column}{0.50\textwidth}

\textbf{Recap --- RQ4:}
\begin{quote}
\small\itshape
How do the three architectures (monolithic, REST, Kafka) compare on
\textbf{latency}, \textbf{throughput}, \textbf{data loss}, and
\textbf{cross-contamination}?
\end{quote}

\vspace{0.5em}
\textbf{This half of the talk answers RQ4 with measured data.}

\vspace{0.5em}
\textbf{What you will see:}
\begin{itemize}
    \item One headline figure per sub-question
    \item Per-application breakdowns (8 apps, light $\rightarrow$ heavy)
    \item Where each architecture wins, and where it breaks
\end{itemize}

\end{column}
\begin{column}{0.46\textwidth}

\begin{block}{Slide Map (figure $\rightarrow$ question)}
\small
\begin{tabular}{@{}ll@{}}
\toprule
\textbf{Figure} & \textbf{Answers} \\
\midrule
fig01 relay        & latency (per app) \\
fig02 vs.\ load    & where it saturates \\
fig03 throughput   & latency under load \\
fig04 save/restore & checkpoint cost \\
fig05 overview     & headline summary \\
fig06 reliability  & loss \& leakage \\
fig07 slowdown     & degradation rate \\
\bottomrule
\end{tabular}
\end{block}

\end{column}
\end{columns}

\end{frame}

% ============================================================
% SLIDE 3: METHODOLOGY
% ============================================================
\begin{frame}{How the Numbers Were Produced}

\begin{columns}[T]
\begin{column}{0.52\textwidth}

\textbf{Load generation: k6}
\begin{itemize}
    \item 10 benchmark scripts (\texttt{00-baseline} $\rightarrow$ \texttt{09-websocket})
    \item Per-app, baseline-calibrated thresholds (\texttt{k6/config.js})
    \item \textbf{5 repeated runs} per architecture for statistical validity
    \item Stepped concurrent-user load (1 $\rightarrow$ 100 VUs)
\end{itemize}

\vspace{0.4em}
\textbf{Deployment: Kubernetes (Rancher)}
\begin{itemize}
    \item Each architecture deployed in isolation
    \item Server-side and remote (laptop) load generators
\end{itemize}

\vspace{0.4em}
\textbf{Extraction \& plotting}
\begin{itemize}
    \item \texttt{k6/extract-results.sh} $\rightarrow$ 3 tidy CSVs
    \item \texttt{analysis/make-plots.R} $\rightarrow$ 7 figures (PNG + PDF)
\end{itemize}

\end{column}
\begin{column}{0.44\textwidth}

\begin{block}{Tidy CSV schema}
\scriptsize
\texttt{summary-long.csv}\\
\quad run, test, metric, avg, med, min, max, p90, p95\\[0.4em]
\texttt{per-app-long.csv}\\
\quad run, test, metric, app\_name, count, avg, med, p95\\[0.4em]
\texttt{per-bucket-long.csv}\\
\quad \dots, vu\_bucket, app\_name, count, avg, med, p95
\end{block}

\begin{alertblock}{Headline statistic}
\small
We report \textbf{p95 latency} (95th percentile), not the mean ---
it reflects the experience of the slowest 1-in-20 interactions, which
is what users actually notice.
\end{alertblock}

\end{column}
\end{columns}

\end{frame}

% ============================================================
% SLIDE 4: TEST SUITE (8 APPS BY WEIGHT)
% ============================================================
\begin{frame}{The Test Suite: 8 Apps, Light to Heavy}

\begin{center}
\small
\begin{tabular}{@{}lll@{}}
\toprule
\textbf{Application} & \textbf{Computational stress} & \textbf{Why it is in the suite} \\
\midrule
Calculator      & baseline (trivial compute) & isolates pure network/relay cost \\
Analytics       & medium data filtering       & typical interactive dashboard \\
Geospatial      & delta-based map sync        & partial-state update path \\
DataExchange    & CSV file relay              & payload-size sensitivity \\
Advanced        & Plotly interactive          & rich-output serialisation \\
ClimateAnomaly  & large CSV + raster          & big-payload stress \\
MonteCarlo      & heavy simulation            & CPU-bound backend \\
MLTrainer       & long-running training       & worst-case compute time \\
\bottomrule
\end{tabular}
\end{center}

\vspace{0.6em}
\begin{center}
\small
Apps are ordered by computational weight throughout the figures, so the
\textbf{left side isolates architecture overhead} and the
\textbf{right side exposes backend saturation}.
\end{center}

\end{frame}

% ============================================================
% SLIDE 5: fig05 OVERVIEW (headline first)
% ============================================================
\begin{frame}{Headline: Latency Overview Across All Benchmarks}

\begin{columns}[T]
\begin{column}{0.62\textwidth}
\begin{center}
\safefig{fig05_overview_grid}{\linewidth}
\end{center}
\end{column}
\begin{column}{0.36\textwidth}
\textbf{What it shows}
\begin{itemize}
    \item One headline p95 per benchmark test
    \item REST vs.\ Kafka side by side
    \item Mean of per-run p95 (5 runs)
\end{itemize}

\vspace{0.5em}
\textbf{Read it as:} the single-slide answer to ``which architecture is
faster, overall?''

\vspace{0.5em}
\begin{block}{Talking point}
\small
Lower bars = lower latency. Look for tests where the two architectures
\textit{diverge} --- those are where the design choice actually matters.
\end{block}
\end{column}
\end{columns}

\end{frame}

% ============================================================
% SLIDE 6: fig01 STATE RELAY LATENCY PER APP
% ============================================================
\begin{frame}{Latency (1/3): State Relay, Per Application}

\begin{columns}[T]
\begin{column}{0.62\textwidth}
\begin{center}
\safefig{fig01_relay_per_app}{\linewidth}
\end{center}
\end{column}
\begin{column}{0.36\textwidth}
\textbf{Metric}
\begin{itemize}
    \item Round-trip: POST $\rightarrow$ backend $\rightarrow$ GET
    \item p95, error bars = SD across runs
\end{itemize}

\vspace{0.5em}
\textbf{Why it matters}\\
This is the \textbf{core collaboration path} --- one user's change
reaching another. It is the operation users perform most.

\vspace{0.5em}
\begin{exampleblock}{Expected story}
\small
On the \textbf{Calculator} (trivial compute) the bar is almost pure
\textit{architecture overhead}: REST polling vs.\ Kafka push shows here
with no compute noise.
\end{exampleblock}
\end{column}
\end{columns}

\end{frame}

% ============================================================
% SLIDE 7: fig02 LATENCY VS CONCURRENT USERS
% ============================================================
\begin{frame}{Latency (2/3): Where Does It Break Under Load?}

\begin{columns}[T]
\begin{column}{0.62\textwidth}
\begin{center}
\safefig{fig02_throughput_vs_load}{\linewidth}
\end{center}
\end{column}
\begin{column}{0.36\textwidth}
\textbf{Metric}
\begin{itemize}
    \item p95 latency vs.\ concurrent users (VUs)
    \item Stepped load on the throughput test
\end{itemize}

\vspace{0.5em}
\textbf{The ``knee''}\\
The point where the curve bends upward marks where the
\textbf{R backend saturates} --- requests start queuing.

\vspace{0.5em}
\begin{block}{Talking point}
\small
A \textbf{later, gentler knee} = better scalability. This is where the
central Spring Boot relay (REST) is expected to bottleneck before
Kafka's partitioned consumers do.
\end{block}
\end{column}
\end{columns}

\end{frame}

% ============================================================
% SLIDE 8: fig03 THROUGHPUT LATENCY PER APP
% ============================================================
\begin{frame}{Latency (3/3): Latency Under Load, Per Application}

\begin{columns}[T]
\begin{column}{0.62\textwidth}
\begin{center}
\safefig{fig03_throughput_per_app}{\linewidth}
\end{center}
\end{column}
\begin{column}{0.36\textwidth}
\textbf{Metric}
\begin{itemize}
    \item p95 latency at 100-VU load, per app
    \item REST vs.\ Kafka
\end{itemize}

\vspace{0.5em}
\textbf{Separating two effects}
\begin{itemize}
    \item \textbf{MLTrainer / MonteCarlo}: compute-bound --- latency is
          mostly R, similar across architectures
    \item \textbf{Lighter apps}: latency reflects \textbf{backend
          queuing} --- the architecture difference dominates
\end{itemize}
\end{column}
\end{columns}

\end{frame}

% ============================================================
% SLIDE 9: fig07 SLOWDOWN FACTOR
% ============================================================
\begin{frame}{Throughput: How Gracefully Does Each App Degrade?}

\begin{columns}[T]
\begin{column}{0.62\textwidth}
\begin{center}
\safefig{fig07_slowdown_factor}{\linewidth}
\end{center}
\end{column}
\begin{column}{0.36\textwidth}
\textbf{Metric}
\begin{itemize}
    \item p95 normalised to each app's lightest load ($=1\times$)
    \item Higher = degrades more under concurrency
\end{itemize}

\vspace{0.5em}
\textbf{Why normalise?}\\
It removes the absolute-speed advantage and asks a fairer question:
\textbf{``how much worse does this app get as users pile on?''}

\vspace{0.5em}
\begin{alertblock}{The money quote}
\small
``At 100 users, app \textit{X} is \textbf{N$\times$ slower} under REST
but only \textbf{M$\times$ slower} under Kafka.''
\end{alertblock}
\end{column}
\end{columns}

\end{frame}

% ============================================================
% SLIDE 10: fig04 SAVE / RESTORE
% ============================================================
\begin{frame}{Persistence: Save \& Restore Checkpoint Cost}

\begin{columns}[T]
\begin{column}{0.62\textwidth}
\begin{center}
\safefig{fig04_save_restore}{\linewidth}
\end{center}
\end{column}
\begin{column}{0.36\textwidth}
\textbf{Metric}
\begin{itemize}
    \item Median latency of the PostgreSQL snapshot cycle
    \item Per application
\end{itemize}

\vspace{0.5em}
\textbf{Architecture-independent}\\
Both REST and Kafka share the \textbf{same DB persistence path}, so this
isolates the cost of the \textit{state-sharing feature itself} (RQ2),
not the sync transport.

\vspace{0.5em}
\begin{block}{Talking point}
\small
Cost scales with \textbf{payload size}: ClimateAnomaly (large raster)
and DataExchange (CSV) sit at the top; Calculator at the bottom.
\end{block}
\end{column}
\end{columns}

\end{frame}

% ============================================================
% SLIDE 11: fig06 RELIABILITY
% ============================================================
\begin{frame}{Reliability: Data Loss, Errors, and Isolation}

\begin{columns}[T]
\begin{column}{0.58\textwidth}
\begin{center}
\safefig{fig06_reliability}{\linewidth}
\end{center}
\end{column}
\begin{column}{0.40\textwidth}
\textbf{Three failure modes, one figure}
\begin{itemize}
    \item \textbf{Data loss} --- did every update arrive?
    \item \textbf{Error rate} --- failed requests under load
    \item \textbf{Cross-contamination} --- did Session A's
          state leak into Session B?
\end{itemize}

\vspace{0.5em}
\begin{alertblock}{The bar that must be zero}
\small
Cross-contamination is a \textbf{correctness}, not a performance,
property. Anything above zero is a defect --- speed cannot buy it back.
\end{alertblock}

\vspace{0.4em}
\small\textit{Lower is better; 0 is ideal across all three.}
\end{column}
\end{columns}

\end{frame}

% ============================================================
% SLIDE 12: MONOLITHIC BASELINE CAVEAT
% ============================================================
\begin{frame}{A Note on the Monolithic Baseline}

\begin{columns}[T]
\begin{column}{0.50\textwidth}

\textbf{Monolithic Shiny is measured differently --- on purpose.}

\vspace{0.5em}
It has \textbf{no REST/Kafka collaboration API}; state moves over
Shiny's internal WebSocket. So k6 cannot trigger its computations the
same way.

\vspace{0.5em}
\textbf{What we still measure with k6:}
\begin{itemize}
    \item HTTP page-load time
    \item WebSocket connection establishment
    \item Max concurrent connections before refusal
\end{itemize}

\end{column}
\begin{column}{0.46\textwidth}

\begin{block}{Computation baseline}
\small
Pure R compute time is captured separately via
\texttt{system.time()} instrumentation inside each monolithic app, then
parsed from container logs.
\end{block}

\vspace{0.5em}
\begin{exampleblock}{Why this is fair}
\small
The monolith is the \textbf{floor} for latency (no network) but the
\textbf{ceiling} for collaboration (it has none). It frames the cost
that REST and Kafka pay \textit{to gain multi-user collaboration}.
\end{exampleblock}

\end{column}
\end{columns}

\end{frame}

% ============================================================
% SLIDE 13: SYNTHESIS / VERDICT
% ============================================================
\begin{frame}{Synthesis: Which Architecture Wins?}

\begin{center}
\small
\begin{tabular}{@{}lccc@{}}
\toprule
\textbf{Dimension} & \textbf{Monolithic} & \textbf{REST + Redis} & \textbf{Kafka} \\
\midrule
Single-user latency      & best (no network) & --- & --- \\
Relay latency (light app)& n/a               & --- & --- \\
Latency under load       & n/a               & --- & --- \\
Scalability (knee)       & connection-limited& --- & --- \\
Reliability / isolation  & single-session    & --- & --- \\
Operational complexity   & lowest            & medium & highest \\
\bottomrule
\end{tabular}
\end{center}

\vspace{0.6em}
\begin{center}
\footnotesize\itshape
Dashes are filled from the measured figures once the 5-run benchmarks
complete. The narrative arc: monolith sets the floor; REST buys
collaboration at modest cost but bottlenecks at the central relay;
Kafka pays higher operational complexity for better behaviour under
concurrency.
\end{center}

\end{frame}

% ============================================================
% SLIDE 14: THREATS TO VALIDITY
% ============================================================
\begin{frame}{Threats to Validity}

\begin{columns}[T]
\begin{column}{0.48\textwidth}
\textbf{Internal}
\begin{itemize}
    \item Shared cluster $\rightarrow$ noisy-neighbour effects;
          mitigated by isolated deployments + 5 repeats
    \item R--Kafka integration maturity may understate Kafka's ceiling
    \item Polling interval (500\,ms) is a tunable that shapes REST latency
\end{itemize}
\end{column}
\begin{column}{0.48\textwidth}
\textbf{External / construct}
\begin{itemize}
    \item 8 apps approximate, but do not exhaust, scientific Shiny usage
    \item Synthetic k6 load $\neq$ real human interaction cadence
    \item p95 chosen as the user-facing statistic; tail beyond p95 not
          the focus
    \item Remote vs.\ server load generators bound network variance
\end{itemize}
\end{column}
\end{columns}

\end{frame}

% ============================================================
% SLIDE 15: CONCLUSION
% ============================================================
\begin{frame}{Conclusion \& Next Steps}

\begin{columns}[T]
\begin{column}{0.52\textwidth}

\textbf{Contribution}
\begin{itemize}
    \item First documented \textbf{collaborative, containerised} R Shiny
          system with two interchangeable sync architectures
    \item A \textbf{reproducible benchmark harness}: k6 + tidy CSVs +
          scripted R figures, runnable end-to-end
\end{itemize}

\vspace{0.5em}
\textbf{Findings (to finalise from data)}
\begin{itemize}
    \item Latency, throughput, reliability quantified per app
    \item Clear REST-vs-Kafka trade-off surface for practitioners
\end{itemize}

\end{column}
\begin{column}{0.44\textwidth}

\begin{block}{Remaining work}
\small
\begin{enumerate}
    \item Complete Kafka 5-run suite (server + remote)
    \item Re-run REST 5$\times$ for statistical validity
    \item Monolithic load + \texttt{system.time()} baseline
    \item Regenerate figures, fill synthesis table
\end{enumerate}
\end{block}

\vspace{0.5em}
\begin{exampleblock}{Future work}
\small
BioDT Honeybee re-engineering; NaaVRE integration; richer concurrency
models beyond write--write.
\end{exampleblock}

\end{column}
\end{columns}

\end{frame}

% ============================================================
% SLIDE 16: THANK YOU
% ============================================================
\begin{frame}{}
\begin{center}
{\Large\textbf{Thank you}}\\[1em]
{\large Questions?}\\[2em]
\small
Iva Gunjaca --- \texttt{ShinySwarm}\\
Master of Software Engineering, University of Amsterdam\\[0.5em]
\footnotesize
Reproduce the figures:\\
\texttt{Rscript analysis/make-plots.R k6/results analysis/plots rest kafka}
\end{center}
\end{frame}

\end{document}
